\tzpanchor{0336}
\tzpentryheader{0336}{Metamaterial--Neural Flow Isomorphism}

\begingroup\small
\begingroup\scriptsize
\setlength{\tabcolsep}{4pt}
\renewcommand{\arraystretch}{0.92}
\noindent\begin{tabularx}{\linewidth}{@{}p{0.24\linewidth}p{0.36\linewidth}X@{}}
\textbf{UUID} & \multicolumn{2}{l@{}}{\tzpcode{9a3a47b8-c5e9-5a63-9836-36cc7caa276f}}\\
\textbf{PiID} & \tzpcode{2540917153643} & \textbf{TZPi}\quad \tzpcode{7}\quad \textbf{PiPE}\quad \tzpcode{343}\\
\textbf{RTEID} & \textbf{Poloidal $\theta$}\quad \tzpcode{1304.7626 rad} & \textbf{Toroidal $\phi$}\quad \tzpcode{02.1498 rad}\\
\textbf{TZPID} & \tzpcode{ID0336} & \textbf{Node Dependencies}\quad \tzpidref{0335}\\
\textbf{Registry} & \tzpcode{registry\_id\_0336} & \textbf{Scale}\quad informational\\
\textbf{LOC Classification} & QA76.9 information and coding & QA614 topological data analysis\\
\textbf{arXiv Classification} & Primary: cs.IT & Cross-list: math.AT, math-ph\\
\textbf{Primary Support} & Metamaterial Avalanche Law & Neural Avalanche Law\\
\textbf{Closure Support} & Isomorphism Constraint & \\
\end{tabularx}\par
\endgroup
\endgroup

\tzpsection{Canonical Statement}
This isomorphism identifies structural and dynamical similarities between metamaterial flow in the TZPQVS system and neural information routing.

\tzpsection{Canonical Equation}
\[
P_{\mathrm{meta}}(s)\sim s^{-\tau}
\]
\[
P_{\mathrm{neuro}}(s)\sim s^{-\tau}
\]

\begin{minipage}[t]{0.49\linewidth}
\tzpsection{Canonical Symbol Key}
\begingroup
\small
\raggedright
\textbf{Source equation symbols.}\par
\begin{itemize}[leftmargin=1.35em,itemsep=1pt,topsep=1pt,parsep=0pt]
    \item $P_{\mathrm{meta}}(s)$: source equation symbol.
    \item $\tau$: source equation symbol.
    \item $P_{\mathrm{neuro}}(s)$: source equation symbol.
    \item $\Phi_{\mathrm{mn}}$: flux, phase, or topological map term in the source equation.
\end{itemize}
\endgroup
\end{minipage}\hfill
\begin{minipage}[t]{0.47\linewidth}
\tzpsection{Wolfram Form}
\begingroup
\scriptsize
\raggedright
\textbf{Source Wolfram model.}\par
\begin{Verbatim}[fontsize=\tiny,breaklines=true,breakanywhere=true]
(* TZPID ID0336 generated source Wolfram model *)
ClearAll[K0336, Cr0336, T, R, A, W, I, N];
eq0336$1 = HoldForm[SimilarTo[PMeta[s], s^(-tau)]];
eq0336$2 = HoldForm[SimilarTo[PNeuro[s], s^(-tau)]];
eq0336$3 = HoldForm[PhiMn[PMeta] == PNeuro];

K0336 = HoldForm[{eq0336$1, eq0336$2, eq0336$3}];

N = "normalized TRAWIN closure object for Metamaterial--Neural Flow Isomorphism";
I = "Isomorphism Constraint integration/comparison layer";
W = "Neural Avalanche Law wave or operator carrier";
A = "Metamaterial Avalanche Law admissible action/amplitude condition";
R = "Neural Avalanche Law rotation/curl preservation layer";
T = "Metamaterial Avalanche Law local source condition";

Cr0336[expr_] := HoldForm[N[I[W[A[R[T[expr]]]]]]];

Result0336 = Cr0336[K0336];
\end{Verbatim}
\endgroup
\end{minipage}

\par\vspace{0.45\baselineskip}
\noindent\begin{minipage}[t]{\linewidth}
\begingroup
\small
\raggedright
\textbf{TRAWIN Operator Key.}\par
\noindent\textbf{Time/localization operator:} $\mathsf{T}\!\left[\mathrm{local}\right]$ Metamaterial Avalanche Law local source condition.\par
\noindent\textbf{Rotation/curl operator:} $\mathsf{R}\!\left[\nabla\times\right]$ Neural Avalanche Law rotation/curl preservation layer.\par
\noindent\textbf{Action/amplitude operator:} $\mathsf{A}\!\left[\mathrm{admissible}\right]$ Metamaterial Avalanche Law admissible action/amplitude condition.\par
\noindent\textbf{Wave/d'Alembertian operator:} $\mathsf{W}\!\left[\Box\right]$ Neural Avalanche Law wave or operator carrier.\par
\noindent\textbf{Integration operator:} $\mathsf{I}\!\left[\int\right]$ Isomorphism Constraint integration/comparison layer.\par
\noindent\textbf{Normalization operator:} $\mathsf{N}\!\left[\mathrm{normalized}\right]$ normalized TRAWIN closure object for Metamaterial--Neural Flow Isomorphism.\par
\endgroup
\end{minipage}

\clearpage
\noindent\textbf{[ID 0336] Metamaterial--Neural Flow Isomorphism}\par
\vspace{0.35em}
\tzpsection{Derivation Layer}
\paragraph{Primary Equation.}
\begin{equation}
\mathcal{V}_T\phi_{\omega}
=
\sqrt{
\frac{\hbar\omega}{2}
}
\,\phi_{\omega}
\coth\left(
\frac{\omega}{2\omega_0}
\right).
\end{equation}

\paragraph{Primary Derivation.}
Let $\phi_{\omega}$ be a field mode of frequency $\omega$.  
The zero-point amplitude of a harmonic mode is proportional to:

\begin{equation}
\sqrt{\frac{\hbar\omega}{2}} .
\end{equation}

In the TZP setting, this amplitude is modulated by a spectral thermal-like factor:

\begin{equation}
\coth\left(
\frac{\omega}{2\omega_0}
\right),
\end{equation}

where $\omega_0$ is the characteristic frequency determined by the spectral density scale.

Acting on the mode $\phi_{\omega}$ gives:

\begin{equation}
\boxed{
\mathcal{V}_T\phi_{\omega}
=
\sqrt{
\frac{\hbar\omega}{2}
}
\,\phi_{\omega}
\coth\left(
\frac{\omega}{2\omega_0}
\right).
}
\end{equation}

\paragraph{Interpretation.}
The vacuum fluctuation operator assigns each mode a zero-point amplitude corrected by the TZP spectral scale.

\par\vspace{0.35em}
\noindent\textbf{Derivable Consequences.}\quad
\begingroup
\footnotesize
\raggedright
The canonical equation anchors Metamaterial--Neural Flow Isomorphism by connecting the source condition to the derivation layer and proof certificate.
\par\endgroup

\tzpsection{Equation Support}
\noindent\textbf{1. Metamaterial Avalanche Law}\par
\[
P_{\mathrm{meta}}(s)\sim s^{-\tau}
\]
\noindent This fixes the avalanche scaling law on the metamaterial side of the proposed correspondence.\par
\noindent\textbf{2. Neural Avalanche Law}\par
\[
P_{\mathrm{neuro}}(s)\sim s^{-\tau}
\]
\noindent This matches it with the homologous neural-avalanche scaling law on the biological side.\par
\noindent\textbf{3. Isomorphism Constraint}\par
\[
\Phi_{\mathrm{mn}}(P_{\mathrm{meta}})=P_{\mathrm{neuro}}
\]
\noindent This closes the entry by stating the isomorphism that carries one critical flow family into the other.\par

\clearpage
\noindent\textbf{[ID 0336] Metamaterial--Neural Flow Isomorphism}\par
\vspace{0.35em}
\tzpsection{Dictionary}
\begingroup\small
\tzpsection{Definition}
Metamaterial--Neural Flow Isomorphism is a registry definition in Field Theory and Action Principles with secondary pressure from Manifold and Metric Dynamics.

% BEGIN TZP CONTEXTUAL MEANING
\tzpsection{Contextual Meaning}
\begingroup\footnotesize\setlength{\emergencystretch}{3em}\sloppy
\textbf{Formal statement:} This isomorphism identifies structural and dynamical similarities between metamaterial flow in the TZPQVS system and neural information routing. \textbf{Contextual meaning:} The entry reads Metamaterial--Neural Flow Isomorphism as a controlled technical headword whose equation layer, proof route, and registry role are interpreted together. \textbf{Derivable consequences:} The entry now reads as an actual critical-scaling correspondence instead of recycling the same quaternary normalization shell. \textbf{Lemma support:} Metamaterial Avalanche Law; Neural Avalanche Law; Isomorphism Constraint.\par
\endgroup
% END TZP CONTEXTUAL MEANING

\tzpsection{Technical Interpretation}
The vacuum fluctuation operator assigns each mode a zero-point amplitude corrected by the TZP spectral scale.

\tzpsection{Equation Grounding}
Equation anchors: 3/2tau approximately 3/2; approx; and Both systems exhibit self-organized criticality with the same avalanche exponent tauapproximately3/2tau approximately 3/2tauapproximately3/2. Proof route: show that the metamaterial and neural avalanche statistics can be matched by an isomorphism that preserves the shared critical exponent. Formalization route: prove that the stated correspondence map preserves the avalanche scaling form between the two carrier systems.

\tzpsection{Interpretive Role}
The entry now reads as an actual critical-scaling correspondence instead of recycling the same quaternary normalization shell.
\endgroup

\tzpsection{TZP Thesaurus}
\begin{enumerate}[leftmargin=1.6em,itemsep=6pt,topsep=4pt]
\item \textbf{Cross-Disciplinary Critical Exponents}: The exact mathematical mapping showing that magnetic flux avalanches and biological brain waves share identical statistical scaling.
\item \textbf{Homologous Phase Transitions}: The direct translation of magnetic susceptibility states into corresponding neural network excitation levels.
\item \textbf{Universal Transport Class Matching}: The proof that widely different physical substrates obey the exact same underlying rules for information propagation.
\end{enumerate}

\tzpsection{Encyclopedia Mathematica}
\begingroup\footnotesize\sloppy
The visual plate below is reserved for the source-specific equation and progression graphic for [ID 0336]. It is included only as a companion to the canonical equation, not as a substitute for the formal text.
\par\endgroup
\ifTZPDarkEdition
\tzpdictionaryvisualband{D:/TZPID/kdp_workflow/build/tikz_figures/ID0336_dictionary_figure_dark.pdf}
\fi

\clearpage
\begingroup\tzpsupportsetup
\noindent\textbf{\fontsize{10.8pt}{12.8pt}\selectfont [ID 0336] Constructive Logic and Proof Certificate}\par
\noindent\textbf{\fontsize{10pt}{12pt}\selectfont Metamaterial--Neural Flow Isomorphism}\par
\vspace{0.45em}
\begingroup
\footnotesize
\setlength{\parskip}{0.22em}
\setlength{\parindent}{0pt}
\noindent\textbf{Curry--Howard--Lambek Interpretation}\par
Under the Curry--Howard--Lambek correspondence, this registry entry is read as a typed proof
object: the stated source condition supplies the proposition, the derivation supplies the
morphism, and the proof certificate records the machine handles needed to check the obligation
in the formal registry.

\vspace{0.45em}
\noindent\textbf{CHL Proof}\par
\textbf{Statement:} this isomorphism identifies structural and dynamical similarities between metamaterial
flow in the TZPQVS system and neural information routing

\textbf{Logic Validation:}\\
\textbf{Proposition:} ``Metamaterial--Neural Flow Isomorphism is valid as a formal dynamical law when the
Hamiltonian or energy operator is well defined on the stated state space.''\\
\textbf{Proof:} The source statement identifies an energy, Hamiltonian, or vacuum-sector mechanism. The
CHL proof reads it as an operator claim: if the state space and localized terms are
defined, then the evolution or extraction channel is formally meaningful.

\textbf{VERDICT:} \ensuremath{\checkmark}\ \textbf{TRUE} --- formal dynamical-operator validation

\textbf{Formal Status:} \ensuremath{\checkmark}\ THEORETICAL -- source-backed CHL proof obligation

\textbf{Physical Status:} THEORETICAL -- requires model-specific empirical interpretation
\par\vspace{0.45em}
\noindent{\tiny\textbf{Isabelle/HOL.} \tzpplaincode{physics\_validity\_from\_model, external\_validity\_package, registry\_number\_matches, equation\_count\_matches}}\par
\ifTZPDarkEdition
\tzpproofvisualband{D:/TZPID/kdp_workflow/build/tikz_figures/ID0336_proof_certificate_figure_dark.pdf}
\fi
\endgroup
\endgroup
